%%%%%%%%%%%%%%%%%%%%%%%%%%%%%%%%%%%%%%%%%
% Masters/Doctoral Thesis 
% LaTeX Template
% Version 2.3 (25/3/16)
%
% This template has been downloaded from:
% http://www.LaTeXTemplates.com
%
% Version 2.x major modifications by:
% Vel (vel@latextemplates.com)
%
% This template is based on a template by:
% Steve Gunn (http://users.ecs.soton.ac.uk/srg/softwaretools/document/templates/)
% Sunil Patel (http://www.sunilpatel.co.uk/thesis-template/)
%
% Template license:
% CC BY-NC-SA 3.0 (http://creativecommons.org/licenses/by-nc-sa/3.0/)
%
%%%%%%%%%%%%%%%%%%%%%%%%%%%%%%%%%%%%%%%%%

%------------------------------------------------------------------------------
%	PACKAGES AND OTHER DOCUMENT CONFIGURATIONS
%------------------------------------------------------------------------------

\documentclass[
11pt, % The default document font size, options: 10pt, 11pt, 12pt
%oneside, % Two side (alternating margins) for binding by default, uncomment to switch to one side
%chapterinoneline,% Have the chapter title next to the number in one single line
%english, % ngerman for German
spanish,
singlespacing, % Single line spacing, alternatives: onehalfspacing or doublespacing
%draft, % Uncomment to enable draft mode (no pictures, no links, overfull hboxes indicated)
%nolistspacing, % If the document is onehalfspacing or doublespacing, uncomment this to set spacing in lists to single
%liststotoc, % Uncomment to add the list of figures/tables/etc to the table of contents
%toctotoc, % Uncomment to add the main table of contents to the table of contents
parskip, % Uncomment to add space between paragraphs
%nohyperref, % Uncomment to not load the hyperref package
headsepline, % Uncomment to get a line under the header
]{MastersDoctoralThesis} % The class file specifying the document structure



\usepackage[utf8]{inputenc} % Required for inputting international characters
\usepackage[T1]{fontenc} % Output font encoding for international characters

\usepackage{palatino} % Use the Palatino font by default
%,style=authoryear
\usepackage[backend=bibtex,natbib=true]{biblatex} % Use the bibtex backend with the authoryear citation style (which resembles APA)

\addbibresource{references.bib} % The filename of the bibliography

\usepackage[autostyle=true]{csquotes} % Required to generate language-dependent quotes in the bibliography

\usepackage{caption}
\usepackage{subcaption}

%------------------------
\usepackage{listings}

\usepackage[framemethod=tikz]{mdframed} % @Eric for License

%\usepackage[hyphens]{url}
%\usepackage[hidelinks]{hyperref}
%\hypersetup{breaklinks=true}
\urlstyle{same}
%\usepackage{cite}

%--------------------------

\usepackage{color}

%
%------------------------------------------------------------------------------
%	MARGIN SETTINGS
%------------------------------------------------------------------------------

\geometry{
	paper=a4paper, % Change to letterpaper for US letter
	inner=2cm, % Inner margin
	outer=3.3cm, % Outer margin
	bindingoffset=2cm, % Binding offset
	top=1.5cm, % Top margin
	bottom=1.5cm, % Bottom margin
	%showframe,% show how the type block is set on the page
}

%------------------------------------------------------------------------------
%	INFORMACIÓN DE LA MEMORIA
%------------------------------------------------------------------------------

\thesistitle{sAPI (simpleAPI): diseño e implementación de una biblioteca para sistematizar
la programación de sistemas embebidos} % El títulos de la memoria, se usa en la carátula y se puede usar el cualquier lugar del documento con el comando \ttitle
\supervisor{MSc. Ing. Félix Gustavo E. Safar.} % El nombre del director, se usa en la carátula y se puede usar el cualquier lugar del documento con el comando \supname
\degree{Magíster en Sistemas Embebidos} % Nombre del grado, se usa en la carátula y se puede usar el cualquier lugar del documento con el comando \degreename
\author{Esp. Ing. Eric Nicolás Pernia} % Tu nombre, se usa en la carátula y se puede usar el cualquier lugar del documento con el comando \authorname
\juradoUNO{Dr. Ing. Pablo M. Gómez (UBA)} % Nombre y pertenencia del un jurado se usa en la carátula y se puede usar el cualquier lugar del documento con el comando \jur1name
\juradoDOS{Mg. Ing. Pablo O. Ridolfi (UTN FRBA)} % Nombre y pertenencia del un jurado se usa en la carátula y se puede usar el cualquier lugar del documento con el comando \jur2name
\juradoTRES{Ing. Juan Manuel Cruz (FIUBA,UTN-FRBA)} % Nombre y pertenencia del un jurado se usa en la carátula y se puede usar el cualquier lugar del documento con el comando \jur3name
\fechaINICIO{marzo de 2018}
\fechaFINAL{diciembre de 2018}

% Your subject area, this is not currently used anywhere in the template, print it elsewhere with \subjectname
\subject{Memoria del Trabajo Final de la Maestría en Sistemas Embebidos (MSE) de la UBA} 
\keywords{MSE, Sistemas Embebidos, CIAA} % Keywords for your thesis, this is not currently used anywhere in the template, print it elsewhere with \keywordnames
\university{Universidad de Buenos Aires} % Your university's name and URL, this is used in the title page and abstract, print it elsewhere with \univname
\faculty{{Facultad de Ingeniería}} % Your faculty's name and URL, this is used in the title page and abstract, print it elsewhere with \facname
\department{Departamento de Electrónica} % Your department's name and URL, this is used in the title page and abstract, print it elsewhere with \deptname
\group{{Laboratorio de Sistemas Embebidos}} % Your research group's name and URL, this is used in the title page, print it elsewhere with \groupname


\hypersetup{pdftitle=\ttitle} % Set the PDF's title to your title
\hypersetup{pdfauthor=\authorname} % Set the PDF's author to your name
\hypersetup{pdfkeywords=\keywordnames} % Set the PDF's keywords to your keywords


\newcaptionname{spanish}{\acknowledgementname}{Agradecimientos}
\newcaptionname{spanish}{\authorshipname}{Declaración de Autoría}
\newcaptionname{spanish}{\abbrevname}{Glosario}
\newcaptionname{spanish}{\byname}{por}

\renewcommand{\lstlistingname}{Algoritmo}% Listing -> Algorithm

% List of Listings -> List of Algorithms
\renewcommand{\lstlistlistingname}{Índice de \lstlistingname s}

\renewcommand{\listtablename}{Índice de Tablas}
\renewcommand{\tablename}{Tabla} 

\addtolength{\footnotesep}{2mm} % Espacio adicional en los footnotes

\begin{document}

% Use roman page numbering style (i, ii, iii, iv...) for the pre-content pages
\frontmatter 

% Default to the plain heading style until the thesis style is called for the
% body content
\pagestyle{plain}

%------------------------------------------------------------------------------
%	CARÁTULA
%------------------------------------------------------------------------------

\begin{titlepage}
\begin{center}

{\scshape\LARGE UNIVERSIDAD DE BUENOS AIRES\par}\vspace{0.1cm} % University name
{\scshape\LARGE FACULTAD DE INGENIERÍA\par}\vspace{0.1cm} % Faculty name
{\scshape\LARGE Maestría en Sistemas Embebidos\par}\vspace{1cm} % Thesis type

\includegraphics[width=.3\textwidth]{./Figures/logoFIUBA.png}
\vspace{0.7cm}

\textsc{\Large Memoria del Trabajo Final}\\[0.5cm] % Thesis type

{\huge \bfseries \ttitle\par}\vspace{0.4cm} % Thesis title

\vspace{0.7cm}
\LARGE\textbf{Autor:\\
\authorname}\\ % Author name

\vspace{1cm}

\large
\vspace{10px}
{Director:} \\
{\supname} % Supervisor name
 
\vspace{0.7cm}
Jurados:\\
\jurunoname\\
\jurdosname\\
\jurtresname
 
\vfill
\textit{Este trabajo fue realizado en las Ciudad Autónoma de Buenos Aires, entre \fechaINICIOname \hspace{1px} y \fechaFINALname.}
\end{center}
\end{titlepage}


%------------------------------------------------------------------------------
%	LICENCIA - LICENSE
%------------------------------------------------------------------------------
\newpage
\begin{center}

\begin{mdframed}[outerlinecolor=black, outerlinewidth=2pt, linecolor=gray, middlelinewidth=3pt, roundcorner=10pt]

 \begin{center}
  \includegraphics*{Figures/cc-by-sa.png}
 \end{center}

\emph{sAPI (simpleAPI): diseño e implementación de una biblioteca para
sistematizar la programación de sistemas embebidos} por Esp. Ing. Eric Nicolás
Pernia se distribuye bajo una \textbf{Licencia Creative Commons
Atribución-CompartirIgual 4.0 Internacional}.
Para ver una copia de esta licencia, visita:
\\
\url{http://creativecommons.org/licenses/by-sa/4.0/}.


\end{mdframed}

\end{center}

%------------------------------------------------------------------------------
%	RESUMEN - ABSTRACT 
%------------------------------------------------------------------------------

\begin{abstract}
\addchaptertocentry{\abstractname} % Add the abstract to the table of contents
%
%The Thesis Abstract is written here (and usually kept to just this page). 
%The page is kept centered vertically so can expand into the blank space above
%the title too\ldots
\centering

En esta memoria se presenta el diseño de una biblioteca para la programación de
sistemas embebidos portable entre plataformas de hardware y lenguajes de
programación. Se realizó una implementación de referencia en lenguaje C para
las plataformas del Proyecto CIAA. Se realizó un banco de pruebas de hardware,
junto a la utilización de testeo unitario e integración contínua para la
validación.

Además de la biblioteca, se creó una herramienta de código abierto para
desarrolladores que automatiza la implemetación de bibliotecas a partir de un
modelo que las describe. De esta manera se facilita la futura ampliación de la
biblioteca y su implementación en otras plataformas de hardware.

\end{abstract}

%------------------------------------------------------------------------------
%	CONTENIDO DE LA MEMORIA  - AGRADECIMIENTOS
%------------------------------------------------------------------------------

\begin{acknowledgements}
% Descomentando esta línea se puede agregar los agradecimientos al índice
%\addchaptertocentry{\acknowledgementname} 
\vspace{1.5cm}

Agradecimientos personales. \textbf{[OPCIONAL]} 

No olvidarse de agradecer al tutor.

No vale poner anti-agradecimientos (este trabajo fue posible a pesar de...)

\end{acknowledgements}

%------------------------------------------------------------------------------
%	LISTA DE CONTENIDOS/FIGURAS/TABLAS
%------------------------------------------------------------------------------
\renewcommand{\listtablename}{Índice de Tablas}

\tableofcontents % Prints the main table of contents

\listoffigures % Prints the list of figures

\listoftables % Prints the list of tables


%------------------------------------------------------------------------------
%	CONTENIDO DE LA MEMORIA  - DEDICATORIA
%------------------------------------------------------------------------------

% escribir acá si se desea una dedicatoria
\dedicatory{\textbf{Dedicado a... [OPCIONAL]}}  

%------------------------------------------------------------------------------
%	CONTENIDO DE LA MEMORIA  - CAPÍTULOS
%------------------------------------------------------------------------------

\mainmatter % Begin numeric (1,2,3...) page numbering

\pagestyle{thesis} % Return the page headers back to the "thesis" style

\renewcommand{\tablename}{Tabla} 

%------------------------------------------------------------------------------

% Define some commands to keep the formatting separated from the content 
\newcommand{\keyword}[1]{\textbf{#1}}
\newcommand{\tabhead}[1]{\textbf{#1}}
\newcommand{\code}[1]{\texttt{#1}}
\newcommand{\file}[1]{\texttt{\bfseries#1}}
\newcommand{\option}[1]{\texttt{\itshape#1}}
\newcommand{\grados}{$^{\circ}$}

%------------------------------------------------------------------------------

%------------------------------------------------------------------------------

% Incluir los capítulos como archivos separados desde la carpeta Chapters
% Descomentar las líneas a medida que se escriben los capítulos

% Introducción General
\chapter{Introducción General}
\label{ChapterIntroGral}
%%------------------------------------------------------------------------------
%	SECTION
%------------------------------------------------------------------------------
\section{Motivación}
\label{sec:motivacion}

En la actualidad existe una enorme variedad de plataformas de sistemas embebidos en el mercado, y si bien todas cuentan con dispositivos programables con características similares y periféricos compatibles, se observa que a la hora de programarlas todas son distintas. En consecuencia, el profesional desarrollador de sistemas embebidos debe invertir mucho tiempo en aprender a programar cada nueva plataforma.

Estas diferencias se deben a que los fabricantes de los dispositivos y empresas asociadas despliegan sus respectivas arquitecturas de hardware tanto en núcleos de procesamiento, como en periféricos. Cada periférico exhibe externamente la funcionalidad esperada, pero su configuración y uso es diferente. Si bien esto da a los fabricantes la posibilidad de diferenciarse frente a la competencia y a los desarrolladores elegir el dispositivo programable que más se adecue a un proyecto, también causa que el programador deba conocer en detalle la arquitectura en particular del dispositivo a la hora de utilizarlo.

Además, estas empresas en su mayoría se limitan a ofrecer información de bajo nivel para programar el hardware directamente (escribiendo registros), o bien, mediante sus propias bibliotecas escritas generalmente en lenguaje C, que están diseñadas con una gran dependencia de la arquitectura de hardware subyacente, con una limitada o nula abstracción del hardware.

Por otro lado, se están portando cada vez más lenguajes de programación a las plataformas de sistemas embebidos, que antes se reservaban para las computadoras de propósito general. Esto se debe a que las nuevas plataformas poseen cada vez más poder de procesamiento y memoria, acompañando la complejidad creciente de las aplicaciones deseadas. En consecuencia se utilizan lenguajes de programación más modernos, en parte por su mayor capacidad de abstracción, y en parte para estandarizar el lenguaje de programación de una aplicación en la que intervienen sistemas embebidos y computadoras de propósito general, siendo un ejemplo típico una aplicación de IoT (del inglés \emph{Internet of Things}), donde intervienen sistemas embebidos distribuidos, infraestructura de red y servidores.

Existen múltiples desarrollos de bibliotecas que logran una abstracción de hardware aceptable. Sin embargo, ninguna de estas se ha adoptado como estándar de facto en la industria y el alcance de las mismas se limita a cubrir plataformas de cierto ecosistema en el marco de una empresa o comunidad y soportan un único lenguaje de programación (en general C o C++).

Este trabajo contribuye a la resolución de esta problemática mediante la sistematización de la implementación de bibliotecas para sistemas embebidos. Para ello se presenta el diseño de una biblioteca para la programación de sistemas embebidos modelada independientemente del lenguaje de programación y arquitectura del hardware. 

Asimismo, se realizaron herramientas de código abierto para desarrolladores, que automatizan gran parte de la implementación de la biblioteca en una nueva plataforma en base a una descripción de la misma, así como la definición de nuevos módulos de bibliotecas. De esta manera se facilita la futura ampliación de la biblioteca y su extensión a diferentes plataformas de hardware.		

Utilizando estas herramientas se realizó una implementación de referencia en lenguaje C para las plataformas del Proyecto CIAA \citep{CIAA}.
%\footnote{Siglas de Computadora Industrial Abierta Argentina.} []. 
%Para su validación, se realizó un banco de pruebas de hardware, junto a la utilización de testeo unitario e integración continua.
Para su validación, se desarrollaron ejemplos funcionales y se utilizaron \emph{tests} unitarios e integración continua.

%------------------------------------------------------------------------------
%	SECTION
%------------------------------------------------------------------------------
\section{Motivación}
\label{sec:motivacion}

En la actualidad existe una enorme variedad de plataformas de sistemas embebidos en el mercado, y si bien todas cuentan con dispositivos programables con características similares y periféricos compatibles, se observa que a la hora de programarlas todas son distintas. En consecuencia, el profesional desarrollador de sistemas embebidos debe invertir mucho tiempo en aprender a programar cada nueva plataforma.

Estas diferencias se deben a que los fabricantes de los dispositivos y empresas asociadas despliegan sus respectivas arquitecturas de hardware tanto en núcleos de procesamiento, como en periféricos. Cada periférico exhibe externamente la funcionalidad esperada, pero su configuración y uso es diferente. Si bien esto da a los fabricantes la posibilidad de diferenciarse frente a la competencia y a los desarrolladores elegir el dispositivo programable que más se adecue a un proyecto, también causa que el programador deba conocer en detalle la arquitectura en particular del dispositivo a la hora de utilizarlo.

Además, estas empresas en su mayoría se limitan a ofrecer información de bajo nivel para programar el hardware directamente (escribiendo registros), o bien, mediante sus propias bibliotecas escritas generalmente en lenguaje C, que están diseñadas con una gran dependencia de la arquitectura de hardware subyacente, con una limitada o nula abstracción del hardware.

Por otro lado, se están portando cada vez más lenguajes de programación a las plataformas de sistemas embebidos, que antes se reservaban para las computadoras de propósito general. Esto se debe a que las nuevas plataformas poseen cada vez más poder de procesamiento y memoria, acompañando la complejidad creciente de las aplicaciones deseadas. En consecuencia se utilizan lenguajes de programación más modernos, en parte por su mayor capacidad de abstracción, y en parte para estandarizar el lenguaje de programación de una aplicación en la que intervienen sistemas embebidos y computadoras de propósito general, siendo un ejemplo típico una aplicación de IoT (del inglés \emph{Internet of Things}), donde intervienen sistemas embebidos distribuidos, infraestructura de red y servidores.

Existen múltiples desarrollos de bibliotecas que logran una abstracción de hardware aceptable. Sin embargo, ninguna de estas se ha adoptado como estándar de facto en la industria y el alcance de las mismas se limita a cubrir plataformas de cierto ecosistema en el marco de una empresa o comunidad y soportan un único lenguaje de programación (en general C o C++).

Este trabajo contribuye a la resolución de esta problemática mediante la sistematización de la implementación de bibliotecas para sistemas embebidos. Para ello se presenta el diseño de una biblioteca para la programación de sistemas embebidos modelada independientemente del lenguaje de programación y arquitectura del hardware. 

Asimismo, se realizaron herramientas de código abierto para desarrolladores, que automatizan gran parte de la implementación de la biblioteca en una nueva plataforma en base a una descripción de la misma, así como la definición de nuevos módulos de bibliotecas. De esta manera se facilita la futura ampliación de la biblioteca y su extensión a diferentes plataformas de hardware.		

Utilizando estas herramientas se realizó una implementación de referencia en lenguaje C para las plataformas del Proyecto CIAA \citep{CIAA}.
%\footnote{Siglas de Computadora Industrial Abierta Argentina.} []. 
%Para su validación, se realizó un banco de pruebas de hardware, junto a la utilización de testeo unitario e integración continua.
Para su validación, se desarrollaron ejemplos funcionales y se utilizaron \emph{tests} unitarios e integración continua.

%------------------------------------------------------------------------------
%	SECTION
%------------------------------------------------------------------------------
\section{Motivación}
\label{sec:motivacion}


%------------------------------------------------------------------------------
\subsection{Lenguajes de programación y Hardware en Sistemas Embebidos}

Saraza...


%------------------------------------------------------------------------------
\subsection{Bibliotecas para microcontroladores ofrecidas por los fabricantes}

Saraza...

%------------------------------------------------------------------------------
\subsection{Proyecto CIAA}

De esta forma, es posible programarlos sin necesidad de conocer detalles sobre
la arquitectura subyacente. Este diseño promueve y permite la programación 
independiente del hardware, reduciendo la complejidad general del desarrollo de
sistemas embebidos. Esta biblioteca se utilizará para programar las diferentes
plataformas que componen el Proyecto Computadora Industrial Abierta Argentina 
(CIAA).


Saraza...



En el mercado se encuentra una gran variedad de plataformas basadas en microcontrolador y aunque todas ellas poseen microcontroladores con características similares y periféricos compatibles, sin embargo, sus bibliotecas son muy diferentes. Esto se debe a que cada fabricante y/o empresas asociadas ofrece sus propias bibliotecas en lenguaje C las cuales están diseñadas fuertemente dependientes de cada arquitectura de cada microcontrolador que estas plataformas contienen.
Existen también muchas bibliotecas que logran una buena abstracción del hardware en varias plataformas, pero ninguna se ha adoptado como estándar general. Esto se debe a múltiples causas, entre ellas:


En el sector de la industria automotriz existe un estándar llamado AUTOSAR[] para la estandariza- ción de la  arquitectura de sistemas electrónicos que aún no ha logrado extenderse a otras industrias y cuya definición de bibliotecas propuesta es muy extensa y compleja de implementar y también para aprender a utilizar.
Bibliotecas de drivers basadas en POSIX[] (que podemos hallar en sistemas con en Linux Embebido como Raspberry Pi[]). Su abstracción ha sido muy útil en la estandarización de drivers para PCs con sistema operativo Unix-compatible. Sin embargo, es tan alejada del hardware físico que en la práctica provoca una muy baja utilización por parte de los profesionales electrónicos y afines que se dedican a la programación de sistemas embebidos.

Programadores hobbistas de sistemas embebidos han impulsado la  estandarización de la programación mediante una biblioteca conocida como Wiring[] disponible para múltiples plataformas (como la popular Arduino[]). Esta biblioteca si bien logra una gran facilidad de uso y rápido aprendizaje contiene algunas imprecisiones técnicas que provoca vicios indeseados en el aprendizaje de programación de microcontroladores. También carece de definición de una API para la utilización del periférico temporizador el cual es muy utilizado en un microcontrolador.

En la actualidad existen muchos fabricantes de microcontroladores que adquieren licencias para la fabricación de microcontroladores con núcleos de procesamiento de arquitecturas Cortex M[] diseñados por la empresa ARM[]. Para los mismos existe una biblioteca estándar llamada CMSIS{} para la programación del núcleo de procesamiento y controlador de interrupciones pero que no se extiende a los perifŕicos donde cada fabricante busca diferenciarse de sus competidores.

Otras empresas muy difundidas en el campo de la enseñanza en la programación como Microchip[] proveen bibliotecas muy dependientes del hardware y generadores automatizados de código que si bien en principio aceleran los tiempos de desarrollo, cuando una aplicación requiere modos más avanzados terminan dificultando la programación pues unas configuraciones pisan a las otras provocando que no funcione.

De los ejemplos anteriores se observa que para la realización de una biblioteca estándar que satisfaga a los diferentes usuarios se debe lograr un balance entre:

\begin{itemize}
\item Extensión de la definición de la API.
\item Dependencia del hardware.
\item Nivel de abstracción.
\item Complejidad de aprendizaje y uso.
\item Periféricos y modos soportados.
\item Escalabilidad.
\end{itemize}



Por estos motivos se decide realizar la definición de una API de una biblioteca estándar para microcontroladores  en lenguaje C que supere todas estas dificultades.

%------------------------------------------------------------------------------
%	SECTION
%------------------------------------------------------------------------------
\section{Objetivos y alcance}
\label{sec:objetivosAlance}

En esta sección se definen los objetivos del presente Trabajo Final (sección \ref{subsec:objetivos}) y el alcance (sección \ref{subsec:alcance}) 

De esta forma, es posible programarlos sin necesidad de conocer detalles sobre
la arquitectura subyacente. Este diseño promueve y permite la programación 
independiente del hardware, reduciendo la complejidad general del desarrollo de
sistemas embebidos. Esta biblioteca se utilizará para programar las diferentes
plataformas que componen el Proyecto Computadora Industrial Abierta Argentina 
(CIAA).

\subsection{Objetivos}
\label{subsec:objetivos}

El objetivo de este proyecto es diseñar e implementar una biblioteca de
software para la programación de sistemas embebidos basados en
microcontroladores con las siguientes características: 

\begin{itemize}
   \item Estar modelada independientemente de los lenguajes de programación.
   \item Definir una interfaz de programación de aplicaciones (API) sencilla que abstraiga los modos de uso más comunes de los periféricos típicos que hallados en cualquier microcontrolador del mercado. 
   \item Ser totalmente portable entre diferentes arquitecturas de hardware sobre donde se ejecuta, manteniendo una API uniforme a lo largo de las mismas\footnote{Debe cumplir la función de capa de abstracción de hardware, o \textit{Hardware Abstraction Layer} (HAL),  en inglés}.
\end{itemize}

Dicha biblioteca se deberá implementar en lenguaje C para las plataformas del
Proyecto CIAA.

\subsection{Alcance}
\label{subsec:alcance}

Este Trabajo Final incluye realización de:

\begin{itemize}
   \item Diseño de la biblioteca. Archivos de descripción de la biblioteca mediante diferentes diagramas y código independiente del lenguaje de programación.
   \item Implementación en lenguaje C de la biblioteca para las plataformas de hardware:
   \begin{itemize}
      \item CIAA-NXP.
      \item EDU-CIAA-NXP.
      \item CIAA-Z3R0.
      \item PicoCIAA.
   \end{itemize}
   \item Manual de instalación de las herramientas para utilizar la biblioteca con las plataformas de hardware citadas.
   \item Manual de referencia de la biblioteca.
   \item Ejemplos de utilización.
\end{itemize}



% Introducción Específica
\chapter{Introducción Específica}
\label{ChapterIntroEsp}

%------------------------------------------------------------------------------
%	SECTION
%------------------------------------------------------------------------------
\section{Plataformas del Proyecto CIAA}
\label{sec:ciaaBoards}



En sus comienzos el proyecto CIAA desarrolla una computadora industrial basada en un microcontrolador NXP LPC4337 (JDB 144) Dual-core Cortex-M4+Cortex-M0 a 204MHz, con 1 MB de memoria Flash y 136 KB de memoria SRAM nombrada CIAA-NXP [ ]. En base a este diseño los integrantes del proyecto CIAA diseñan una versión educativa sin las protecciones industriales, nombrada EDU-CIAA-NXP [ ]. En la figura [ ] se presenta esta plataforma.

Esta plataforma incluye todos los periféricos típicos que podemos encontrar en los microcontroladores disponibles en el mercado permitiendo la enseñanza con herramientas modernas.

Mediante la colaboración de la Red Universitaria de Sistemas Embebidos (RUSE) [ ] se han distribuido en 2015 entre 10 y 40 placas en Universidades de Argentina con carreras afines a la electrónica.

Utilizando la EDU-CIAA-NXP como plataforma base para la enseñanza se ha desarrollado, o colaborado en el desarrollo de diferentes herramientas de software y hardware abiertas, las cuales se describen en las siguientes secciones.

--------------------------------------


El proyecto de la Computadora Industrial Abierta Argentina (CIAA) nació en 2013 como una iniciativa conjunta entre el sector académico y el industrial, representados por la ACSE\footnote{Asociación Civil para la investigación, promoción y desarrollo de los Sistemas electrónicos Embebidos. Sitio web: \url{http://www.sase.com.ar/asociacion-civil-sistemas-embebidos}} y CADIEEL\footnote{Cámara Argentina de Industrias Electrónicas, Electromecánicas y Luminotécnicas. Sitio web: \url{http://www.cadieel.org.ar/}}, respectivamente.

\medskip

\noindent Los objetivos del proyecto CIAA son:

\begin{itemize}
\item
Impulsar el desarrollo tecnológico nacional, a partir de sumar valor agregado al trabajo y a los productos y servicios, mediante el uso de sistemas electrónicos, en el marco de la vinculación de las instituciones educativas y el sistema científico-tecnológico con la industria.
\item
Darle visibilidad positiva a la electrónica argentina.
\item
Generar cambios estructurales en la forma en la que se desarrollan y utilizan en nuestro país los conocimientos en el ámbito de la electrónica y de las instituciones y empresas que hacen uso de ella.
\end{itemize}

\noindent Todo esto en el marco de un trabajo libre, colaborativo y articulado entre industria y academia. 

\bigskip

Con esta iniciativa, se han desarrollado en la actualidad varias plataformas de hardware y entornos de programación para utilizarlas.

\medskip

Al momento de la presentación de este trabajo, existen dos versiones de la plataforma CIAA cuyo desarrollo ha sido completado:

\begin{itemize}
\item 
CIAA-NXP, basada en el microcontrolador NXP LPC4337, que ya se comercializa.
\item 
CIAA-FSL, que utiliza, en cambio, el microcontrolador Freescale MK60FX512VLQ15, pero únicamente hay prototipos de esta plataforma. 
\end{itemize}

Además, existe una versión educativa de bajo costo de la CIAA-NXP, nombrada EDU-CIAA-NXP, que ya se distribuyeron alrededor de 1000 unidades y ya hay otras 1000 reservadas en producción. 

\medskip

Debido a estas razones, el trabajo se enfoca en el desarrollo de herramientas para programar las dos plataformas basadas en el microcontrolador NXP LPC4337. Se introducen a continuación las características de las mismas.




Siendo el autor participante de este proyecto desde mediados de 2014, ocupando el rol de Responsable de Software-PLC mediante el aporte al proyecto CIAA de un IDE\footnote{IDE4PLC. Sitio web: \url{http://proyecto-ciaa.com.ar/devwiki/doku.php?id=desarrollo:software-plc}} que permite programar esta plataforma con lenguajes de PLC industriales (IEC-661131-3), se desea agregar en esta oportunidad la posibilidad de programar a esta plataforma con un lenguaje de programación orientado a objetos mediante el desarrollo de un IDE para tal fin.


\subsection{CIAA-NXP}

La CIAA-NXP es la primera y única computadora del mundo que reúne dos cualidades:

\begin{itemize}
\item 
Ser \textbf{Industrial}, ya que su diseño está preparado para las exigencias de confiabilidad, temperatura, vibraciones, ruido electromagnético, tensiones, cortocircuitos, etc., que demandan los productos y procesos industriales.
\item 
Ser \textbf{Abierta}, ya que toda la información sobre su diseño de hardware, firmware, software, etc. está libremente disponible en Internet bajo la Licencia BSD, para que cualquiera la utilice como quiera.
\end{itemize}

\noindent Esta plataforma se compone de:

\begin{itemize}
\item
CPU: Microcontrolador NXP LPC 4337 JDB 144 (Dual-core Cortex-M4 + Cortex-M0 @ 204MHz).
\item
Debugger: USB-to-JTAG FT2232H. Soportado por OpenOCD.
\item
Memorias: 
   \begin{itemize}
   \item
   IS42S16400F - SDRAM. 64Mbit @ 143MHz.
   \item
   S25FL032P0XMFI011 - Flash SPI. 32 Mbit, Quad I/O Fast read: 80 MHz.
   \item
   24AA1025 - EEPROM I2C. 1 Mbit, 400 kHz. Almacenamiento de propósito general, datos de calibración del usuario, etc.
   \item
   24AA025E48 - EEPROM I2C. 2 kbit, 400 kHz. Para implementación de MAC-Address o almacenamiento de propósito general.
   \end{itemize}
\item
Entradas y salidas:   
   \begin{itemize}
   \item
   8 entradas digitales opto-aisladas 24VDC.
   \item
   4 Entradas analógicas 0-10V/4-20mA.
   \item
   4 salidas Open-Drain 24VDC.
   \item
   4 Salidas con Relay DPDT.
   \item
   1 Salida analógica 0-10V/4-20mA.
   \end{itemize}
\item
LV-GPIO:
   \begin{itemize}
   \item
   14 GPIOs.
   \item
   I2C.
   \item
   SPI.
   \item
   4 canales analógicos.
   \item
   Aux. USB.
   \end{itemize}
\item
Interfaces de comunicación:
   \begin{itemize}
   \item
   Ethernet.
   \item
   USB On-The-Go.
   \item
   RS232.
   \item
   RS485.
   \item
   CAN.
   \end{itemize}
\item
Múltiples fuentes de alimentación.

\end{itemize}

\noindent En la figura [\ref{fig:ciaaNxp}] se muestra una fotografía de la plataforma.

\begin{figure}[!htbp]
\begin{center}  %[width=14cm,height=8cm]
\includegraphics*[width=\textwidth]{Figures/CIAA-NXP_Foto.png}
\par\caption{Plataforma CIAA-NXP.}\label{fig:ciaaNxp}
\end{center}
\end{figure}


%http://proyecto-ciaa.com.ar/devwiki/doku.php?id=desarrollo:hardware:ciaa_nxp:ciaa_nxp_inicio

\subsection{EDU-CIAA-NXP}

La plataforma EDU-CIAA-NXP es un desarrollo colaborativo, realizado por miembros de la Red Universitaria de Sistemas Embebidos (RUSE), en el marco del Proyecto CIAA. RUSE se compone de docentes pertenecientes a más de 60 Universidades a lo largo y a lo ancho del país. 

\medskip

\noindent Los propósitos de la plataforma son: 

\begin{itemize}
\item
Proveer una plataforma de desarrollo moderna, económica y de fabricación nacional basada en la CIAA-NXP, que sirva a docentes y a estudiantes en los cursos de sistemas embebidos.
\item
Lograr una amplia inserción en el sistema educativo argentino.
\item
Realizar un aporte eficaz al desarrollo de vocaciones tempranas en electrónica, computación e informática.
\item
Demostrar que las universidades argentinas son capaces de realizar un desarrollo colaborativo exitoso en el área de los sistemas embebidos, cumpliendo con requerimientos de tiempo y forma.
\end{itemize}

\medskip

\noindent Características de la EDU-CIAA-NXP:

\begin{itemize}
\item
CPU: Microcontrolador NXP LPC 4337 JDB 144 (Dual-core Cortex-M4 + Cortex-M0 @ 204MHz).
\item
Debugger: USB-to-JTAG FT2232H. Soportado por OpenOCD.
\item
2 puertos micro-USB (uno para aplicaciones y debug, otro OTG).
\item
6 salidas digitales implementadas con leds (3 normales y uno RGB).
\item
4 entradas digitales con pulsadores.
\item
1 puerto de comunicaciones RS-485 con bornera.
\item
2 conectores de expasión:
   \begin{itemize}
   \item
   P0:
      \begin{itemize}
      \item
      3 entradas analógicas (ADC0 a ADC2).
      \item
      1 salida analógica (DAC0).
      \item
      1 conexión para un teclado de 3 x 4.
      \item
      12 pines genéricos de I/O.
      \end{itemize}
   \item
   P1:
      \begin{itemize}
      \item
      1 puerto Ethernet.
      \item
      1 puerto CAN.
      \item
      1 puerto SPI.
      \item
      1 puerto I2C.
      \item
      12 pines genéricos de I/O.
      \end{itemize}
   \end{itemize}
\end{itemize}

\noindent En la figura [\ref{fig:eduCiaa}] se muestra una fotografía de esta plataforma.

\begin{figure}[!htbp]
\begin{center}  %[width=14cm,height=8cm]
\includegraphics*[width=\textwidth]{Figures/EDU-CIAA-NXP_Foto.png}
\par\caption{Plataforma EDU-CIAA-NXP.}\label{fig:eduCiaa}
\end{center}
\end{figure}

\medskip 

%http://proyecto-ciaa.com.ar/devwiki/doku.php?id=desarrollo:edu-ciaa:edu-ciaa-nxp

% http://proyecto-ciaa.com.ar/devwiki/doku.php?id=start

\subsection{Pico-CIAA}

http://www.proyecto-ciaa.com.ar/devwiki/doku.php?id=desarrollo:hardware:picociaa

Procesador Dual-Core de 32 bits

Tiene un procesador Dual Core de 32-bit ARM Cortex-M4F / Cortex-M0+ @100 MHz con 512kB flash y 104 kB SRAM, incluye debug vía un LPC11U35 programable.

Diagrama en bloques

Posee diversos puertos de comunicación, entre ellos USB, PCIExpress, UART, SPI, I2C y soporte para PWM y entradas y salidas digitales de propósito general.
special

Es la CIAA más pequeña

Con sólo 51 x 30 mm la picoCIAA es la placa más pequeña de la familia CIAA,
lo que la hace ideal para aplicaciones de internet de las cosas (IoT).

Ideal en Single Board Computers

Su interfaz PCI Express permite integrarla fácilmente en SBC de propósito general, para dar soporte de bajo nivel al manejo de dispositivos en tiempo real.

\begin{figure}[!htbp]
\begin{center}  %[width=14cm,height=8cm]
\includegraphics*[width=\textwidth]{Figures/PicoCIAA_Foto.png}
\par\caption{Plataforma PicoCIAA.}\label{fig:picoCiaa}
\end{center}
\end{figure}

\subsection{CIAA-Z3R0}

Esta plataforma es ideal para aplicaciones de bajo consumo y proyectos de robótica educativa. Está diseñada para ser integrada como componente en un diseño mayor debido a su reducido tamaño conectándola mediante tiras de pines, o soldada a través de su borde de agujeros para montaje castellated.

Se puede comprar en Argentina por aproximadamente AR \$400 a inicios de 2018.

Posee la mayoría de los periféricos que se encuentran en la EDU-CIAA-NXP pero utiliza un microcontrolador siete veces más económico que esta última, de la empresa Silicon Labs, modelo EFM32HG322F64 (QFP48) con núcleo ARM Cortex-M0+ a 25MHz, 64KB de memoria Flash y 8KB de memoria SRAM, que es suficiente para que un alumno entre en el mundo de los microcontroladores modernos de 32 bits.

El dispositivo de depuración se debe comprar por separado, sin embargo, mediante un único dispositivo de depuración se pueden programar muchas plataformas CIAAZ3R0.
No es necesario dejar este circuito de depuración en el diseño final

\begin{figure}[!htbp]
\begin{center}  %[width=14cm,height=8cm]
\includegraphics*[width=\textwidth]{Figures/CIAA-Z3R0_Foto.png}
\par\caption{Plataforma CIAA-Z3R0.}\label{fig:ciaaZero}
\end{center}
\end{figure}

\subsection{Arquitectura de Hardware}

Saraza...

\subsection{Bibliotecas disponibles}

Saraza...
 
%------------------------------------------------------------------------------
%	SECTION
%------------------------------------------------------------------------------
\section{Requerimientos}
\label{sec:requerimientos}

%\subsection{Diseño de la biblioteca}

%\subsection{Implementación}

%\subsection{Documentación y difusión}

Los requerimientos se establecieron en base a los objetivos expuestos en la sección \ref{sec:objetivosAlance}, reuniones con desarrolladores del Proyecto CIAA, alumnos (de grado de UNQ, posgrado de FIUBA y cursos CAPSE) y el director. El espíritu de los mismos es proponer una solución que permita la estandarización de la programación de microcontroladores. Estos son:

\begin{enumerate}
   \item Fecha de finalización: 19/11/2018.
   \item Diseño de la biblioteca.
      \begin{enumerate}[1]
         \item Realizar un diseño independiente del hardware y lenguaje de programación, teniendo en cuenta los conceptos familiares al programador de Sistemas Embebidos.
         \item Debe estar modelada con objetos y contar con una descripción mediante diagramas UML.
         \item Debe modelar al menos:
            \begin{enumerate}[1]
               \item CORE: un núcleo de procesamiento.
               \item GPIO: periférico que consiste en un único pin de entrada/salida de propósito general (pin), así como un grupo de pines (port).
               \item ADC: periférico conversor analógico-digital.
               \item DAC: periférico conversor digital-analógico.
               \item TIMER: periférico temporizador.CORE: un núcleo de procesamiento.
               \item RTC: periférico reloj de tiempo real.CORE: un núcleo de procesamiento.
               \item UART: periférico de comunicación serial asincrónico.
               \item SPI: periférico interfaz serie sincrónica.
               \item I2C: periférico de comunicación serie entre circuitos integrados.
            \end{enumerate}
      \end{enumerate}
   \item Implementación de la biblioteca.
      \begin{enumerate}[1]
         \item Utilizar un sistema de control de versiones con repositorios on line.
         \item Programar en lenguaje C la biblioteca para cada plataforma de hardware particular utilizando como plantilla los archivos generados.
         \item Desarrollar ejemplos de utilización para las diferentes plataformas.
      \end{enumerate}
   \item Documentación y difusión.
      \begin{enumerate}[1]
         \item Confeccionar un manual de referencia de la API de la biblioteca.
         \item Desarrollar un tutorial de instalación de las herramientas para utilizar la biblioteca con las plataformas de hardware citadas.
         \item Publicación on line del código fuente.
         \item Informar a la comunidad del Proyecto CIAA y a la comunidad de programadores de
   Sistemas Embebidos.
      \end{enumerate}
\end{enumerate}

% \subsection{Entregables principales del trabajo}
% 
% Se listan a continuación los entregables principales del proyecto:
% 
% 1. Archivos de descripción de la biblioteca mediante diferentes diagramas y código
% independiente del lenguaje de programación.
% 2. Código fuente de la implementación en lenguaje C de la biblioteca para las plataformas de hardware:
%    - CIAA-NXP.
%    - EDU-CIAA-NXP.
%    - CIAA-Z3R0.
%    - PicoCIAA.
% 3. Ejemplos de utilización.
% 4. Informe final.
% 5. Manual de referencia de la API de la biblioteca.
% 6. Tutorial de instalación de las herramientas para utilizar la biblioteca con las plataformas de hardware citadas.

 
%------------------------------------------------------------------------------
%	SECTION
%------------------------------------------------------------------------------
\section{Planificación}
\label{sec:planificacion}

En la tabla \ref{tab:planificacion} se resume la planificación del trabajo. En la misma se pueden observar cada una de las tareas planeadas, junto a su duración, fecha de inicio y fecha de finalización estimadas.

\begin{table}[!htbp]
\centering
\includegraphics*[width=16cm]{Figures/GanttSP.png}
\par\caption{Planificación del Trabajo Final.}\label{tab:planificacion}
\end{table} 

% Diseño
\chapter{Diseño}
\label{ChapterDiseno}

%------------------------------------------------------------------------------
%	SECTION
%------------------------------------------------------------------------------
\section{Descripción general}
\label{sec:descripGralDiseno}

El sistema propuesto para la creación de bibliotecas de software para la programación de plataformas de hardware se ilustra en la figura \ref{fig:sapi_gen}.

\begin{figure}[!htbp]
\begin{center}  % [width=14cm,height=8cm] [width=\textwidth]
\includegraphics*[width=14cm]{Figures/sapi_gen.pdf}
\par\caption{Diagrama del sistema diseñado.}\label{fig:sapi_gen}
\end{center}
\end{figure}

Se diseñó un modelo de plataforma de hardware (en adelante \emph{board} por sus siglas en inglés). Este modelo describe una plataforma de hardware completa y se debe instanciar para cada nueva plataforma de hardware.

Por otra parte, se desarrollaron módulos de biblioteca (en adelante módulos sAPI) independientes del hardware y lenguajes de programación. Estos definen propiedades y métodos describen las propiedades y métodos de cada uno de los periféricos de un cierto SoC, definiendo entonces, la API de la biblioteca.

Además, se debe proveer la implementación de cada uno de los métodos de los módulos sAPI (en adelante \emph{drivers}). Estos \emph{drivers} deben estar escritos en un cierto lenguaje de programación y son dependientes de la arquitectura del hardware.

También, se diseñó un generador de biblioteca sAPI. Este generador se debe definir para cada lenguaje de programación. Para este trabajo final se desarrolló solamente el generador de lenguaje C, sin embargo, el mismo se realizó de forma modular para que pueda ser fácilmente adaptado a otros lenguajes. 

Mediante el generador de C, utilizando una instancia de \emph{board} y \emph{drivers} específicos, junto con los módulos sAPI, se genera entonces una biblioteca sAPI en lenguaje C para una plataforma particular.

En las siguientes secciones de detalla el diseño e implementación de cada una de las entidades descriptas. Las mismas se modelaron utilizaron los conceptos del paradigma de la programación orientada a objetos, que son especialmente útiles para describir módulos de software que encapsulan funcionalidad. Estos se describen mediante diagramas UML de clases de forma independiente del lenguaje de programación.


%------------------------------------------------------------------------------
%	SECTION
%------------------------------------------------------------------------------
\section{Modelado de plataforma de hardware genérica}
\label{sec:modelHardware}

%------------------------------------------------------------------------------
\subsection{Clasificación de módulos de hardware}


%------------------------------------------------------------------------------
\subsection{Plataforma y componentes}


%------------------------------------------------------------------------------
\subsection{SoC}

%------------------------------------------------------------------------------
\subsection{Módulo Core}

Events:
Reset
%void ResetISR(void);
%WEAK void NMI_Handler(void);
%WEAK void HardFault_Handler(void);
%WEAK void MemManage_Handler(void);
%WEAK void BusFault_Handler(void);
%WEAK void UsageFault_Handler(void);
%WEAK void SVC_Handler(void);
%WEAK void DebugMon_Handler(void);
%WEAK void PendSV_Handler(void);
%WEAK void SysTick_Handler(void);
%WEAK void IntDefaultHandler(void);


sleep(UntilNextInterrupt)?
sleep(UntilGPIOPinIsHigh)

%------------------------------------------------------------------------------
\subsection{Módulo SoC}


%------------------------------------------------------------------------------
\subsection{Módulo Board}



%void M0CORE_IRQHandler(void) ALIAS(IntDefaultHandler);
%
%void SCT_IRQHandler(void) ALIAS(IntDefaultHandler);
%void RIT_IRQHandler(void) ALIAS(IntDefaultHandler);
%
%void TIMER0_IRQHandler(void) ALIAS(IntDefaultHandler);
%void TIMER1_IRQHandler(void) ALIAS(IntDefaultHandler);
%void TIMER2_IRQHandler(void) ALIAS(IntDefaultHandler);
%void TIMER3_IRQHandler(void) ALIAS(IntDefaultHandler);
%
%void ADC0_IRQHandler(void) ALIAS(IntDefaultHandler);
%
%void I2C0_IRQHandler(void) ALIAS(IntDefaultHandler);
%
%void SPI_IRQHandler(void) ALIAS(IntDefaultHandler);
%
%void UART0_IRQHandler(void) ALIAS(IntDefaultHandler);
%void UART1_IRQHandler(void) ALIAS(IntDefaultHandler);
%void UART2_IRQHandler(void) ALIAS(IntDefaultHandler);
%void UART3_IRQHandler(void) ALIAS(IntDefaultHandler);
%
%void GPIO0_IRQHandler(void) ALIAS(IntDefaultHandler);
%void GPIO1_IRQHandler(void) ALIAS(IntDefaultHandler);
%void GPIO2_IRQHandler(void) ALIAS(IntDefaultHandler);
%void GPIO3_IRQHandler(void) ALIAS(IntDefaultHandler);
%void GPIO4_IRQHandler(void) ALIAS(IntDefaultHandler);
%void GPIO5_IRQHandler(void) ALIAS(IntDefaultHandler);
%void GPIO6_IRQHandler(void) ALIAS(IntDefaultHandler);
%void GPIO7_IRQHandler(void) ALIAS(IntDefaultHandler);
%void GINT0_IRQHandler(void) ALIAS(IntDefaultHandler);
%void GINT1_IRQHandler(void) ALIAS(IntDefaultHandler);



%------------------------------------------------------------------------------
\subsection{Módulo GPIO}

El módulo GPIO modela tanto un único pin de entrada/salida de propósito general (pin), así como a un conjunto de pines (puerto). 

\titulo{Tipos de datos}

\begin{itemize}
\item
\emph{Pin\_t}, modela un pin.
\item
\emph{Port\_t}, representa un puerto.
\item
\emph{PortValue\_t}, representa el valor de todos los pines del puerto.
\end{itemize}

La cantidad de pines que forman un puerto depende de la arquitectura.

\titulo{Propiedades de pin}

\begin{itemize}
\item
\emph{mode}, del tipo \emph{PinMode\_t}, cuyos valores posibles se listan en la tabla \ref{tab:PinModeT}.


\begin{table}[h]
	\centering	
	\begin{tabular}{l l}   
		\toprule
		\textbf{Valor} 	    & \textbf{Descripción}  \\
		\midrule
		DISABLE	 & Deshabilitado \\		
		INPUT	    & Pin configurado como entrada	\\
		OUTPUT	 & Pin configurado como salida, modo \emph{push-pull}	\\
		OUTPUT\_PUSHPULL & Pin configurado como salida, modo \emph{push-pull}	\\
		OUTPUT\_OPEN\_DRAIN	 & Pin como salida, en modo drenador abierto	\\
		\bottomrule
		\hline
	\end{tabular}
	\caption[Valores posibles del tipo \emph{PinMode\_t}.]{Valores posibles del tipo \emph{PinMode\_t} y su descripción.}
	\label{tab:PinModeT}
\end{table}

\item
\emph{pull}, del tipo \emph{PinPull\_t}. Los valores posibles se listan en la tabla \ref{tab:PullT}.


\begin{table}[h]
	\centering	
	\begin{tabular}{l l}   
		\toprule
		\textbf{Valor} 	    & \textbf{Descripción} \\
		\midrule
		NONE	       & Modo por defecto, pin flotante \\		
		PULL\_UP	    & Pin con resistencia de \emph{pull-up} \\
		PULL\_DOWN	 & Pin con resistencia de \emph{pull-down} \\
		PULL\_BOTH	 & Pin con ambas resistencias \\
		\bottomrule
		\hline
	\end{tabular}
	\caption[Valores posibles del tipo \emph{PinPull\_t}.]{Valores posibles del tipo \emph{PinPull\_t} y su descripción.}
	\label{tab:PullT}
\end{table}

\item
\emph{value}, del tipo \emph{bool\_t}. Representa el valor a escribir o leer de un pin.
\end{itemize}

\titulo{Métodos de pin}

Los métodos básicos de configuración y uso de un pin son:

\begin{verbatim}
   • init( PintInit_t init )
   • read() : bool_t
   • write( bool_t value )
   • deInit()
\end{verbatim}

%Speed. %Specifies the speed for the selected pins. %Possible values are:
%GPIO_SPEED_LOW
%GPIO_SPEED_MEDIUM
%GPIO_SPEED_FAST
%GPIO_SPEED_HIGH

El parámetro \emph{init} se debe completar con un valor posible de \emph{mode}, y opcionalmente, un valor de \emph{pull} separados con $|$, por ejemplo:

\begin{verbatim}
   • init( INPUT )
   • init( INPUT | PULL_UP )
\end{verbatim}

Contiene además el método \emph{toggle()} intercambia el valor del pin si el mismo está configurado como salida.:

\begin{verbatim}
   • toggle()
\end{verbatim}

Para establecer una interrupción en un pin ante cierto evento define los métodos: 

\begin{verbatim}
   • eventCallbackSet( PinEvent_t evt, Callback_t func )
   • eventCallbackClear( PinEvent_t evt )
   • interrupt( bool_t enable ) : PinStatus_t
 \end{verbatim} 

\emph{eventCallbackSet()} establece el evento ante el cual se interrumpe, junto a una función de \emph{callback} que se ejecutará ante su ocurrencia. \emph{eventCallbackClear()} permite remover el \emph{callback}. En la tabla \ref{tab:PinEvtT} se listan los eventos del tipo \emph{PinEvent\_t}.

\begin{table}[h]
	\centering	
	\begin{tabular}{l l}   
		\toprule
		\textbf{Valor} & \textbf{Descripción} \\
		\midrule
		HIGH\_LEVEL  & Interrupción por nivel alto \\		
		LOW\_LEVEL   & Interrupción por nivel bajo \\
		RISING\_EDGE & Interrupción por flanco ascendente \\
		FALLING\_EDGE & Interrupción por flanco descendente \\
		BOTH\_EDGES & Interrupción por flanco ascendente y descendente \\
		\bottomrule
		\hline
	\end{tabular}
	\caption[Valores posibles del tipo \emph{PinEvent\_t}.]{Valores posibles del tipo \emph{PinEvent\_t} y su descripción.}
	\label{tab:PinEvtT}
\end{table}

\emph{interrupt()} permite habilitar o deshabilitar la interrupción de hardware del pin ante la ocurrencia del evento. Según la plataforma particular, existe la posibilidad que soporte interrupción en algunos pines fijos, o que exista una cierta cantidad de interrupciones de pin y se pueda elegir a qué pines se le asignan, es por eso que el método \emph{interrupt()} retorna un valor del tipo \emph{PinStatus\_t}, cuyos valores aceptados se exponen en la tabla \ref{tab:PinStatusT}. 

\begin{table}[h]
	\centering	
	\begin{tabular}{l l}   
		\toprule
		\textbf{Valor} & \textbf{Descripción} \\
		\midrule
		STATUS\_OK & Se pudo configurar la interrupción en el pin \\		
		ERROR\_NOT\_SUPPORT & El pin no soporta capacidad de interrupción \\
		ERROR\_NO\_MORE\_AVAILABLE & Si ya se configuraron todas las interrupciones\\
		 & de hardware disponibles \\
		\bottomrule
		\hline
	\end{tabular}
	\caption[Valores posibles del tipo \emph{PinStatus\_t}.]{Valores posibles del tipo \emph{PinStatus\_t} y su descripción.}
	\label{tab:PinStatusT}
\end{table}

\titulo{Propiedades y métodos de puerto}

Puerto define la propiedad \emph{value}, del tipo \emph{PortValue\_t}. Representa el valor a escribir o leer en un puerto.

Para la configuración y uso de un puerto se definen los métodos:

\begin{verbatim}
   • init( List_t config )
   • read() : PortValue_t
   • write( PortValue_t value )
   • deInit()
\end{verbatim}

El método \emph{init()} de puerto requiere una lista de configuraciones. Los métodos \emph{read()} y \emph{write()} permiten leer y escribir el puerto completo.

Se agregan además los métodos:

\begin{verbatim}
   • pins(): List_t
   • set( PortValue_t mask )
   • reset( PortValue_t mask )
\end{verbatim}

El método \emph{pins()}, devuelve una lista de pines que componen el puerto. Los métodos \emph{set()} y \emph{reset()} permiten escribir el puerto afectado por una máscara de bits que se le pasa como parámetro.

% Para implementacion ver:

% https://github.com/epernia/sAPI/blob/develop/sapi1Project/libs/sapi_soc_peripherals/inc/sapi_gpio.h

% https://github.com/epernia/sAPI/blob/develop/sapi1Project/libs/sapi_soc_peripherals/src/sapi_gpio.c

% https://github.com/epernia/sAPI/blob/develop/toMergeIn_sapi1Project/sapi_soc_peripherals/PERIPH/GPIO/gpioInterruptLPCOpenExample.c.txt

% Ejemplos de micropython de uart interrupt de gigliotti


%------------------------------------------------------------------------------
\subsection{Módulo ADC}

El módulo ADC modela un periférico Conversor Analógico-Digital. Este periférico contiene varios canales de conversión multiplexados los cuales se pueden configurar individualmente como entradas analógicas.

\titulo{Propiedades de ADC}

\begin{itemize}
\item
\emph{conversionRate}, del tipo \emph{int\_t}, representa la tasa de conversión en Hertz. Los valores posibles dependen de cada plataforma. Se pueden utilizar alternativamente los valores genéricos: \emph{LOW}, \emph{MEDIUM}, \emph{HIGH} y \emph{VERY\_HIGH}.
\item
\emph{conversionMode}, del tipo \emph{AdcConvMode\_t}, con valores: \emph{SINGLE} (activado por software desde el programa de usuario), \emph{CONTINUOUS} (conversion periódica disparada por hardware a tasa \emph{conversionRate}).% y \emph{EXTERNAL\_TRIGGERED} (Timer Match signal o GPIO).
\item
\emph{voltageReferece}, del tipo \emph{AdcVRef\_t} cuyos valores posibles son: \emph{VCC}, \emph{INTERNAL} y \emph{EXTERNAL}.
\item
\emph{resolution}, del tipo \emph{AdcRes\_t}, cuyos valores también dependen de la plataforma pero se pueden utilizar en su lugar los valores genéricos.
\item
\emph{channelsMode}, del tipo \emph{AdcChannelsMode\_t}, con valores: \emph{SIGLE} y \emph{DIFFERENTIAL}.
\item
\emph{location}, del tipo \emph{Location\_t}. Son las posibles ubicaciones del periférico con respecto a pines físicos del SoC. Si bien dependen de la arquitectura, se definen los valores genéricos \emph{LOCATION0} a \emph{LOCATION7} para cada plataforma.
\end{itemize}

\titulo{Métodos de ADC}

Inicialización y restablecimiento de un periférico ADC:

\begin{verbatim}
   • init()
   • init( int_t conversionRate, AdcInit_t init )
   • deInit()
\end{verbatim}

Los valores posibles del parámetro \emph{init} se forman con valores de:

\begin{verbatim}
   • conversionMode | voltageReferece | resolution | 
     channelsMode | location
\end{verbatim}

Todos los valores de configuración del método \emph{init()} son opcionales y los que no se aplican se inicializan por defecto con:

\begin{verbatim}
   • conversionRate = VERY_HIGH
   • voltageReferece = VCC
   • conversionMode = SINGLE
   • resolution = VERY_HIGH
   • channelMode = SINGLE
   • loacation = LOCATION0
\end{verbatim}

Habilitación/deshabilitación de un canal particular del ADC (configuración de una entrada analógica):

\begin{verbatim}
   • channel( AdcChannel_t channel, bool_t enable )
\end{verbatim}

Para conocer los canales disponibles de un ADC existe el método 
\emph{channelsGet()} que devuelve la lista de canales del ADC. Esta lista se reduce a la mitad si el modo \emph{channelsMode} es \emph{DIFFERENTIAL)}

\begin{verbatim}
   • channelsGet() : List_t
\end{verbatim}

Si el ADC tiene configurado \emph{conversionMode = CONTINUOUS} muestrea automáticamente los canales habilitados. El método\emph{startConversion()} comienza el muestreo automático, para parar el muestreo se utiliza \emph{stopConversion()}. Si en cambio \emph{conversionMode = SINGLE}, la lectura se debe lanzar mediante el método \emph{startConversion()} que recibe un canal como parámetro. 
Una conversión se puede abortar con \emph{stopConversion()}.

\begin{verbatim}
   • startConversion( AdcChannel_t channel )
   • stopConversion()
\end{verbatim}

Existen 2 formas de leer los valores convertidos del ADC, por encuesta, o por interrupción. En el primer caso se debe chequear si finalizó la conversión con el método \emph{status()}, que retorna un un valor del tipo \emph{AdcStatus\_t}. Si se realizó la conversión correctamente se puede leer el valor convertido con el método \emph{read()}.

\begin{verbatim}
   • status() : AdcStatus_t
   • read() : AdcValue_t
\end{verbatim}

Los valores posibles de \emph{AdcStatus\_t} son: 

%\emph{CONVERTING}, \emph{CONVERSION\_COMPLETE}, \emph{CONVERSION\_COMPLETE\_OVERRUN}, \emph{CONVERSION\_ADC\_ERROR}, \emph{CONVERSION- \_TIMEOUT}. 

\begin{verbatim}
   • READY
   • BUSY
   • CONVERSION_COMPLETE
   • ERROR
   • ERROR_TIMEOUT
\end{verbatim}

Para leer un canal por interrupción se debe primero establecer una función de \emph{callback} al evento \emph{CONVERSION\_COMPLETE}, otra al evento \emph{ERROR} y luego activar la interrupción del periférico. Esto se realiza con los métodos:

\begin{verbatim}
   • eventCallbackSet( ChannelEvent_t evt, Callback_t func )
   • eventCallbackClear( ChannelEvent_t evt )
   • interrupt( bool_t enable )
\end{verbatim}

Los eventos del ADC son del tipo \emph{AdcEvent\_t}. Al completar la conversión se ejecutará una de las dos funciones.

\titulo{Métodos de Canal de ADC}

Como alternativa, existen métodos que aplican a cierto canal de ADC y proveen un mayor nivel de abstracción:

\begin{verbatim}
   • read() : AdcValue_t
   • readSync( AdcValue_t data, Time_t timeout )
               : AdcStatus_t
   • readAsync( Callback_t sucess, Callback_t error )
\end{verbatim}

Tanto \emph{read()} como \emph{readSync()} realizan una lectura bloqueante, en consecuencia, es necesario establecer un tiempo de \textit{time out} para que no bloquee indefinidamente. \emph{read()} fija automáticamente el \emph{timeout} al doble del \emph{conversionRate}, siempre retorna un valor, que a veces puede ser inválido. \emph{readSync()} devuelve el estado de conversión, para que el usuario chequee si valor convertido es válido, en ese caso lo carga en la variable \emph{data}, que recibe como parámetro. Si se establece en 0 el valor de \emph{timeout} en el método con sufijo \emph{Sync} entonces el tiempo de \texttt{timeout} será considerado infinito. Lo mismo sucede con todos los métodos que cuenten con dicho sufijo que se describen que se describen en esta sección.

El método \emph{readAsync()} realiza una lectura no bloqueante. Recibe como parámetros dos funciones de \emph{callback}, una que lanza en caso de conversión exitosa (si esto ocurre le pasa como parámetro el valor convertido) y otra en caso de error (le pasa el estado de conversión).

Este módulo no define método \emph{write()}.

% Ver de agregar Analog Comparator como evento: Comparación del valor convertido contra un valor programado, para mayor que, igual que o menor que.
%    ADC_GREATER_THAN: Modo comparación por mayor que.
%    ADC_LESS_THAN: Modo comparación por menor que.

%----------------------------------------------------------------------------------------
%	SECTION
%----------------------------------------------------------------------------------------
\section{Análisis del software}
 
La idea de esta sección es resaltar los problemas encontrados, los criterios utilizados y la justificación de las decisiones que se hayan tomado.

Se puede agregar código o pseudocódigo dentro de un entorno lstlisting con el siguiente código:

\begin{verbatim}
\begin{lstlisting}[caption= "un epígrafe descriptivo"]

	las líneas de código irían aquí...
	
\end{lstlisting}
\end{verbatim}

A modo de ejemplo:

\begin{lstlisting}[caption=Pseudocódigo del lazo principal de control.]  % Start your code-block

#define MAX_SENSOR_NUMBER 3
#define MAX_ALARM_NUMBER  6
#define MAX_ACTUATOR_NUMBER 6

uint32_t sensorValue[MAX_SENSOR_NUMBER];		
FunctionalState alarmControl[MAX_ALARM_NUMBER];	//ENABLE or DISABLE
state_t alarmState[MAX_ALARM_NUMBER];						//ON or OFF
state_t actuatorState[MAX_ACTUATOR_NUMBER];			//ON or OFF

void vControl() {

	initGlobalVariables();
	
	period = 500 ms;
		
	while(1) {

		ticks = xTaskGetTickCount();
		
		updateSensors();
		
		updateAlarms();
		
		controlActuators();
		
		vTaskDelayUntil(&ticks, period);
	}
}
\end{lstlisting}




%------------------------------------------------------------------------------
\subsection{Módulo Timer}
Duración: TIME.

323 Mazidi

TIMER: modela un periférico Timer/Counter.

Time/Counterl Capture Event OOOC
Time/Counterl Compare Match A OOOE
Time/Counterl Compare Match B 001 0
Time/Counterl Overflow 001 2


\emph{Tick\_t}, representa un valor de conteo de un lapso de tiempo adimencional lapso de tiempo en ticks (se almacena internamente como \emph{uint64\_t}).

Como evento posible se define: \emph{CONVERSION\_COMPLETE} que genera una interrupción cuando se completa la conversión. El método \emph{interrupt()} habilita  o deshabilita todas las interrupciones del periférico.

% Para implementacion ver:

% https://github.com/epernia/sAPI/blob/develop/toMergeIn_sapi1Project/sapi_soc_peripherals/inc/sapi_timer.h

% la nueva sapi en la carpeta dev


%De notas.txt:
%
%
%Para TIMER en general:
%----------------------
%
%    Timer(4, prescaler=624, period=13439, mode=UP, div=1)
%
%The information means that this timer is set to run at the peripheral
%clock speed divided by 624+1, 
%
%it will count from 0 up to 13439, at which point it triggers an interrupt, 
%and then starts counting again from 0.
%
%These numbers are set to make the timer trigger at 10 Hz: the source frequency
%of the timer is 84MHz (found by running ``tim.source_freq()``) so we
%get 84MHz / 625 / 13440 = 10Hz.
%
%Timer counter
%-------------
%
%So what can we do with our timer?  The most basic thing is to get the
%current value of its counter::
%
%    >>> tim.counter()
%    21504
%
%This counter will continuously change, and counts up.
%
%---------------------------------------------------------
%
%
%Para TIMER PWM:
%---------------
%
%	enum POLARITY{ ACTIVE_LOW=0, ACTIVE_HIGH=1 };
%
%	float analogFrequency;  //defaults to 100,000 Hz
%	float analogMax;        //defaults to 3.3V
%
%	int setPeriod(unsigned int period_ns);
%	 unsigned int getPeriod();
%	 int setFrequency(float frequency_hz);
%	 float getFrequency();
%	 int setDutyCycle(unsigned int duration_ns);
%	 int setDutyCycle(float percentage);
%	 unsigned int getDutyCycle();
%	 float getDutyCyclePercent();
%
%	 int setPolarity(PWM::POLARITY);
%	 void invertPolarity();
%	 PWM::POLARITY getPolarity();
%
%	virtual void setAnalogFrequency(float frequency_hz) { this->analogFrequency = frequency_hz; }
%	virtual int calibrateAnalogMax(float analogMax); //must be between 3.2 and 3.4
%	virtual int analogWrite(float voltage);
%
%	virtual int run();
%	virtual bool isRunning();
%	virtual int stop();
%
%	virtual ~PWM();
%private:
%	float period_nsToFrequency(unsigned int);
%	unsigned int frequencyToPeriod_ns(float);
%};
%
%---------------------------------------------------------
%
%Para TIMER Waveform generation output signal:
%----------------------------------------------
%
%   // Generate a periodic sine wave in an array of 100 values - using ints
%   unsigned int waveform[100];
%   float gain = 50.0f;
%   float phase = 0.0f;
%   float bias = 50.0f;
%   float freq = 2.0f * 3.14159f / 100.0f;
%   for (i=0; i<100; i++){
%      waveform[i] = (unsigned int)(bias + (gain * sin((i * freq) + phase)));
%   }
%   
%   
%   
%tickConfig( TICK_MS(1) ); o tickStart( TICK_MS(1) );
%tickEnabeling( ENABLE );
%tickSetCallback( function );
%
%typedef tick_t uint64_t;
%
%#define TICK_MS(t) ((typedef)(t))
%#define TICK_S(t)  (((typedef)(t))*1000)




%------------------------------------------------------------------------------
\subsection{Módulo RTC}

El módulo RTC modela un periférico Reloj de Tiempo Real. Este periférico permite mantener la hora y fecha en un sistema.

\titulo{Propiedades de RTC}

\begin{itemize}
\item
\emph{dateAndTime}, contiene la fecha/hora actual, del tipo \emph{DateAndTime\_t}.
\item
\emph{alarm}, establece la fecha/hora para un evento de alarma, del tipo \emph{DateAndTime\_t}.
\item
\emph{interval}, del tipo \emph{Time\_t}, establece el intervalo de tiempo para un evento periódico.
\end{itemize}

\titulo{Métodos de RTC}

Inicialización y restablecimiento de un periférico RTC:

\begin{verbatim}
   • init( DateAndTime_t absTime | )
   • deinit()
\end{verbatim}

El parámetro \emph{absTime} establece la fecha/hora actual. La lectura del RTC se realiza con el método \emph{read()}, que retorna la fecha/hora actual. Para cambiar la fecha/hora se utiliza el método \emph{write()}:

\begin{verbatim}
   • read() : DateAndTime_t
   • write( DateAndTime_t absTime )
\end{verbatim}

Para establecer una interrupción en un RTC ante cierto evento define los métodos: 

\begin{verbatim}
   • interrupt( bool_t enable )
   • eventCallbackSet( ChannelEvent_t evt, Callback_t c )
   • eventCallbackClear( ChannelEvent_t evt )
\end{verbatim}

Los eventos posibles son: \emph{ALARM}, que genera una única interrupción cuando se cumple cierto valor absoluto de fecha/hora, y \emph{PERIODIC}, que interrumpe de forma periódica cada cierto lapso de tiempo.

%------------------------------------------------------------------------------
\subsection{Módulo UART}

El módulo UART modela un periférico de comunicación serie transmisor/receptor asincrónico universal.

\titulo{Propiedades de UART}

\begin{itemize}
\item
\emph{baudRate}, del tipo \emph{UartDataBits\_t}.
\item
\emph{dataBits}, del tipo \emph{UartDataBits\_t}.
\item
\emph{parity}, del tipo \emph{UartParity\_t}.
\item
\emph{stopBits}, del tipo \emph{UartStopBits\_t}.
\item
\emph{flowControl}, del tipo \emph{UartFlowCtrl\_t}.
\end{itemize}

En la tabla \ref{tab:uartProperties} se listan los posibles valores de configuración de estas propiedades. 

\begin{table}[h]
	\centering	
	\begin{tabular}{l l l l l}   
		\toprule
		\textbf{baudRate} & \textbf{dataBits} & \textbf{parity} & \textbf{stopBits} & \textbf{flowControl} \\
		\midrule
		\textbf{int\_t} & \textbf{UartDataBits\_t} & \textbf{UartParity\_t} & \textbf{UartStopBits\_t} & \textbf{UartFlowCtrl\_t} \\
		\midrule
      1200   & DATABITS5 & NONE & STOPBITS1          & NONE     \\
      2400   & DATABITS6 & EVEN & STOPBITS2          & CTS      \\
      4800   & DATABITS7 & ODD  & STOPBITS1\_5 (1.5) & RTS      \\
      9600   & DATABITS8 &      &                    & RTS\_CTS \\
      19200  & DATABITS9 &      &                    &          \\
      38400  &           &      &                    &          \\
      57600  &           &      &                    &          \\
      115200 &           &      &                    &          \\
		\bottomrule
		\hline
	\end{tabular}
	\caption[Valores posibles de las propiedades de UART]{Valores posibles de las propiedades de UART}
	\label{tab:uartProperties}
\end{table}

\begin{itemize}
\item
\emph{enableTransmitter}, del tipo \emph{bool\_t}, se utiliza para habilitar/deshabilitar el transmisor de la UART.
\item
\emph{enableReceiver}, del tipo \emph{bool\_t}, permite habilitar/deshabilitar el receptor de la UART.
\item
\emph{location}, del tipo \emph{Location\_t}.
\item
\emph{transmitValue}, del tipo \emph{byte\_t}. Representa el valor a transmitir.
\item
\emph{receiveValue}, del tipo \emph{byte\_t}. Representa el último valor recibido. 
\end{itemize}

\titulo{Métodos de UART}

Inicialización y restablecimiento de un periférico UART:

\begin{verbatim}
   • init()
   • init( int_t baudRate )
   • init( int_t baudRate, UartInit init )
   • deInit()
\end{verbatim}

Los valores posibles del parámetro \emph{init} se forman con valores de:

\begin{verbatim}
   • dataBits | parity | stopBits | flowControl |
     enableTransmitter | enableReceiver | location
\end{verbatim}

Todos los valores del método \emph{init()} son opcionales. En caso de no inicializarlos los valores por defecto son:

\begin{verbatim}
   • baudRate = 9600
   • dataBits = DATABITS8
   • parity = NONE
   • stopBits = STOPBITS1
   • flowControl = NONE
   • loacation = LOCATION0
   • enableTransmitter = TRUE
   • enableReceiver = TRUE
\end{verbatim}

Se puede transmitir o recibir un byte usando el módulo UART tanto por encuesta como por interrupción mediante los métodos:

\begin{verbatim}
   • status() : UartStatus_t
   • send( const byte_t value )
   • receive() : byte_t
   • eventCallbackSet( UartEvent_t evt, Callback_t func )
   • eventCallbackClear( UartEvent_t evt )
   • interrupt( bool_t enable )
\end{verbatim}
  
Los posibles valores de \emph{UartStatus\_t} son:

%\emph{TRANSMIT\_COMPLETE}, \emph{RECEIVE\_COMPLETE}, \emph{TRANSMIT\_DATA\_OVERRUN}, \emph{RECEIVE\_FRAME\_ERROR}, \emph{RECEIVE\_DATA\_OVERRUN},  \emph{RECEIVE\_PARITY\_ERROR}, \emph{TRANSMIT\_TIMEOUT} y \emph{RECEIVE\_TIMEOUT}

\begin{verbatim}
   • TRANSMITTER_READY
   • TRANSMIT_COMPLETE
   • RECEIVE_COMPLETE
   • ERROR_TRANSMIT_DATA_OVERRUN
   • ERROR_RECEIVE_FRAMER
   • ERROR_RECEIVE_DATA_OVERRUN
   • ERROR_RECEIVE_PARITY
   • ERROR_TIMEOUT
\end{verbatim}

Eventos (tipo \emph{UartEvent\_t}):

%\emph{TRANSMIT\_COMPLETE}, \emph{RECEIVE\_COMPLETE}, \emph{TRANSMIT\_ERROR}, \emph{RECEIVE\_ERROR} 

\begin{verbatim}
   • TRANSMITTER_READY
   • TRANSMIT_COMPLETE
   • RECEIVE_COMPLETE
   • ERROR_TRANSMIT
   • ERROR_RECEIVE
\end{verbatim}

Existe además el método \emph{sendBreak()} que envía una \emph{break condition} al bus. Esto significa que mantiene la línea de transmisión en nivel bajo durante un tiempo mayor al requerido para enviar un caracter.

\begin{verbatim}
   • sendBreak()
\end{verbatim}

\titulo{Métodos de UART de alto nivel}

Métodos de lectura de byte:

\begin{verbatim}
   • readByte() : byte_t
   • readByteSync( inOut byte_t value, Time_t timout )
                   : UartStatus_t
   • readByteAsync( Callback_t sucess, Callback_t error )
\end{verbatim}

Métodos de escritura de byte:

\begin{verbatim}
   • writeByte( const byte_t value )
   • writeByteSync( const byte_t value, Time_t timout )
                    : UartStatus_t
   • writeByteAsync( const byte_t value, 
                     Callback_t sucess,
                     Callback_t error 
                   )
\end{verbatim}

Métodos de lectura de \emph{String}:

\begin{verbatim}
   • readString( inOut String_t data,
                 inOut size_t dataSize,
                 String_t terminator
               )
   • readStringSync( inOut String_t data,
                     inOut size_t dataSize,
                     String_t terminator,
                     Time_t timout 
                   ) : UartStatus_t
   • readStringAsync( inOut String_t data,
                      inOut size_t dataSize,
   	                  String_t terminator,
                      Callback_t sucess, 
                      Callback_t error 
                    ) 
\end{verbatim}

Métodos de escritura de \emph{String}:                  
                
\begin{verbatim}  
   • writeString( const String_t data
                  inOut size_t dataSize
                )
   • writeStringSync( const String_t data,
                      inOut size_t dataSize,
                      Time_t timeout
                    ) : UartStatus_t
   • writeStringAsync( const String_t value, 
                      inOut size_t dataSize,
                      Callback_t sucess,
                      Callback_t error
                     )
\end{verbatim}

Métodos de lectura de arreglo de bytes:

\begin{verbatim}
   • readByteArray( inOut ByteArray_t data,
                    inOut size_t dataSize
                  )
\end{verbatim}

\pagebreak

\begin{verbatim}
   • readByteArraySync( inOut ByteArray_t data,
                        inOut size_t dataSize,
                        Time_t timeout
                      ) : UartStatus_t
   • readByteArrayAsync( inOut ByteArray_t data, 
                         inOut size_t dataSize,
                         Callback_t sucess,
                         Callback_t error 
                       )
\end{verbatim}

Métodos de escritura de arreglo de bytes:

\begin{verbatim}
  • writeByteArray( const ByteArray_t data,
                    inOut size_t dataSize
                  )
  • writeByteArraySync( const ByteArray_t data,
                        inOut size_t dataSize,
                        Time_t timeout
                      ) : UartStatus_t
  • writeByteArrayAsync( const ByteArray_t data, 
                         inOut size_t dataSize,
                         Callback_t sucess,
                         Callback_t error
                       )
\end{verbatim}

En todos los métodos de lectura se le pasa como primer parámetro la variable donde escribirá los datos recibidos. En la lectura de \emph{String} se le debe pasar un parámetro \texttt{terminator} que es un \texttt{String} que indica como debe ser el fin de cadena a buscar. Cuando se lee un \emph{String} o \emph{Byte Array} se almacena la cantidad de datos leídos en la variable \emph{dataSize} que se le pasa como parámetro.

%------------------------------------------------------------------------------
\subsection{Módulo SPI}

Modela un bus de serie para interfaz con periféricos (bus SPI). El periférico SPI proporciona comunicación serial sincrónica \emph{full duplex} entre dispositivos maestros y esclavos. Se usa comúnmente para la comunicación con periféricos externos como memorias flash, sensores, relojes en tiempo real (RTC), etc.

\titulo{Propiedades de SPI}

\begin{itemize}
\item
\emph{mode}, del tipo \emph{SpiMode\_t}, que representa el modo de funcionamiento del periférico SPI, con valores \emph{MASTER} o \emph{SLAVE}.
\item
\emph{clockFrequency}, del tipo \emph{int\_t}, es la frecuencia de reloj (en Hz) del BUS SPI (solo tiene sentido en modo \emph{MASTER}).
\item
\emph{clockPolarity}, del tipo \emph{SpiClockParity\_t}, define el nivel cuando se considera activo el clock, con valores \emph{POLARITY\_LOW} o \emph{POLARITY\_HIGH}.
\item
\emph{clockPhase}, del tipo \emph{SpiClockPhase\_t}, establece en que flanco de reloj se muestrean los datos, cuyos valores posibles son \emph{PHASE\_EDGE1} o \emph{PHASE\_EDGE2}.
\item
\emph{dataBits}, del tipo \emph{SpiDataBits\_t}, que puede tomar los valores \emph{DATABITS8} o \emph{DATABITS16}.
\item
\emph{dataOrder}, del tipo \emph{SpiDataOrder\_t}, representa el orden con el que se mueven los bits en el registro de desplazamiento, es decir si se mueven primero el bit menos significativo o el más significativo: \emph{MSB\_FIRST} o \emph{LSB\_FIRST}.
\item
\emph{location}, del tipo \emph{Location\_t}.
\item
\emph{transmitValue}, del tipo \emph{ByteArray\_t}. Representa el valor a transmitir.
\item
\emph{receiveValue}, del tipo \emph{ByteArray\_t}. Representa el último valor recibido. 
\end{itemize}

\titulo{Métodos de SPI}

Inicialización y restablecimiento de un periférico SPI:

\begin{verbatim}
   • init()
   • init( SpiInit_t init )
   • init( SpiInit_t init, int_t clockFrequency )
   • deInit()
\end{verbatim}

Los valores posibles del parámetro \emph{init} se forman con valores de:

\begin{verbatim}
   • mode | clockPolarity | clockPhase | dataOrder |
     dataBits | location
\end{verbatim}

Los valores de configuración son opcionales, con valores por defecto:

\begin{verbatim}
   • mode = MASTER
   • clockFrequency = 100000
   • clockPolarity = POLARITY_HIGH
   • clockPhase = PHASE_EDGE1
   • dataBits = DATABITS8
   • dataOrder MSB_FIRST
   • location = LOCATION0
\end{verbatim}

En caso de seleccionar modo \emph{MASTER} se debe configurar la frecuencia de reloj y al menos un pin como selector de esclavo. La frecuencia de reloj puede configurarse usando el método \emph{init()} de 2 parámetros, o bien, mediante:

\begin{verbatim}
   • clockFrequencySet( int_t clockFrequency )
\end{verbatim}

Utilizando \emph{initSSPin()} se configura un pin como selector de esclavo:

\begin{verbatim}
   • initSSPin( Pin_t pin )
\end{verbatim}

Luego para elegir que esclavo habilitar o deshabilitar se utiliza:

\begin{verbatim}
   • slaveSelect( Pin_t pin, bool_t enable )
\end{verbatim}

Permite realizar envío, recepción o transferencia de datos (que envía y recibe simultáneamente). Estas operaciones pueden realizarse por encuesta o interrupción con los siguientes métodos:

\begin{verbatim}
   • status() : SpiStatus_t
   • send( const ByteArray_t data )
   • receive( inOut ByteArray_t receive )
   • transfer( const ByteArray_t send, 
               inOut ByteArray_t receive
             ) : SpiStatus_t
   • eventCallbackSet( ChannelEvent_t evt, Callback_t func )
   • eventCallbackClear( ChannelEvent_t evt )
   • interrupt( bool_t enable )
\end{verbatim}

Los posibles valores de \emph{SpiStatus\_t} son:

\begin{verbatim}
   • READY
   • BUSY
   • TRANSFER_COMPLETE
   • ERROR
   • ERROR_TIMEOUT
\end{verbatim}

Eventos:

\begin{verbatim}
   • TRANSFER_COMPLETE
   • TRANSFER_ERROR
\end{verbatim}

\titulo{Métodos de SPI de alto nivel}

% Los métodos de escritura y lectura de a un byte no tienen sentido pues si el tamaño de datos no es 8, entonces siempre es byte array.

Lectura de arreglo de bytes:

\begin{verbatim}
   • readByteArray( inOut ByteArray_t data,
                    inOut size_t dataSize
                  )
   • readByteArraySync( inOut ByteArray_t data,
                        inOut size_t dataSize,
                        Time_t timout
                      ) : SpiStatus_t
   • readByteArrayAsync( inOut ByteArray_t data, 
                         inOut size_t dataSize,
                         Callback_t sucess,
                         Callback_t error 
                       )
\end{verbatim}

Escritura de arreglo de bytes:

\begin{verbatim}
   • writeByteArray( const ByteArray_t data,
                     inOut size_t dataSize
                   )
   • writeByteArraySync( const ByteArray_t data,
                         inOut size_t dataSize,
                         Time_t timout
                       ) : SpiStatus_t
   • writeByteArrayAsync( const ByteArray_t data, 
                          inOut size_t dataSize,
                          Callback_t sucess,
                          Callback_t error 
                        )
\end{verbatim}

Transferencia de arreglo de bytes:

\begin{verbatim}
   • transferByteArray( const ByteArray_t send,
                        inOut size_t sendSize,
                        inOut ByteArray_t receive,
                        inOut size_t dataSize
                      )
   • transferByteArraySync( const ByteArray_t send,
                            inOut size_t sendSize,
                            inOut ByteArray_t receive,
                            inOut size_t dataSize
                          ) : SpiStatus_t
   • transferByteArrayAsync( const ByteArray_t send,
                             inOut size_t sendSize,
                             inOut ByteArray_t receive,
                             Callback_t sucess,
                             Callback_t error 
                           )
\end{verbatim}

% Para la implementación:

% https://github.com/epernia/sAPI/blob/develop/toMergeIn_sapi1Project/sapi_soc_peripherals/inc/sapi_spi.h

%  https://github.com/epernia/sAPI/blob/develop/toMergeIn_sapi1Project/sapi_soc_peripherals/src/sapi_spi.c

%https://www.mikroe.com/ebooks/programming-dspic-microcontrollers-in-pascal/master-and-slave-modes

%------------------------------------------------------------------------------
\subsection{Módulo I2C}

I2C: modela un periférico de comunicación serie I2C.

Events:
I2C


%------------------------------------------------------------------------------
%	SECTION
%------------------------------------------------------------------------------
\section{Verificación del modelo}
\label{sec:modelVerification}

%------------------------------------------------------------------------------
%\subsection{}



%------------------------------------------------------------------------------
%	SECTION
%------------------------------------------------------------------------------
\section{Modelo abstracto de módulo de biblioteca}
\label{sec:modelLibrary}

Cada módulo que compone la biblioteca sAPI se modela utilizando los conceptos del paradigma de la programación orientada a objetos, los cuales son especialmente útiles para describir módulos de software que encapsulan funcionalidad. Estos se describen mediante diagramas UML de clases de forma independiente del lenguaje de programación.



%----------------------------------------------------------------------------------------
%	SECTION
%----------------------------------------------------------------------------------------
\section{Análisis del software}
 
La idea de esta sección es resaltar los problemas encontrados, los criterios utilizados y la justificación de las decisiones que se hayan tomado.

Se puede agregar código o pseudocódigo dentro de un entorno lstlisting con el siguiente código:

\begin{verbatim}
\begin{lstlisting}[caption= "un epígrafe descriptivo"]

	las líneas de código irían aquí...
	
\end{lstlisting}
\end{verbatim}

A modo de ejemplo:

\begin{lstlisting}[caption=Pseudocódigo del lazo principal de control.]  % Start your code-block

#define MAX_SENSOR_NUMBER 3
#define MAX_ALARM_NUMBER  6
#define MAX_ACTUATOR_NUMBER 6

uint32_t sensorValue[MAX_SENSOR_NUMBER];		
FunctionalState alarmControl[MAX_ALARM_NUMBER];	//ENABLE or DISABLE
state_t alarmState[MAX_ALARM_NUMBER];						//ON or OFF
state_t actuatorState[MAX_ACTUATOR_NUMBER];			//ON or OFF

void vControl() {

	initGlobalVariables();
	
	period = 500 ms;
		
	while(1) {

		ticks = xTaskGetTickCount();
		
		updateSensors();
		
		updateAlarms();
		
		controlActuators();
		
		vTaskDelayUntil(&ticks, period);
	}
}
\end{lstlisting}





% Implementación
\chapter{Implementación}
\label{ChapterImplementacion}

%------------------------------------------------------------------------------
%	SECTION
%------------------------------------------------------------------------------
\section{Descripción general} % Sobre la implementación.
\label{sec:implemDescripGral}

%------------------------------------------------------------------------------
%	SECTION
%------------------------------------------------------------------------------
\section{Archivos de descripción de las paltaformas del proyecto CIAA}


%----------------------------------------------------------------------------------------
%	SECTION
%----------------------------------------------------------------------------------------
\section{Análisis del software}
 
La idea de esta sección es resaltar los problemas encontrados, los criterios utilizados y la justificación de las decisiones que se hayan tomado.

Se puede agregar código o pseudocódigo dentro de un entorno lstlisting con el siguiente código:

\begin{verbatim}
\begin{lstlisting}[caption= "un epígrafe descriptivo"]

	las líneas de código irían aquí...
	
\end{lstlisting}
\end{verbatim}

A modo de ejemplo:

\begin{lstlisting}[caption=Pseudocódigo del lazo principal de control.]  % Start your code-block

#define MAX_SENSOR_NUMBER 3
#define MAX_ALARM_NUMBER  6
#define MAX_ACTUATOR_NUMBER 6

uint32_t sensorValue[MAX_SENSOR_NUMBER];		
FunctionalState alarmControl[MAX_ALARM_NUMBER];	//ENABLE or DISABLE
state_t alarmState[MAX_ALARM_NUMBER];						//ON or OFF
state_t actuatorState[MAX_ACTUATOR_NUMBER];			//ON or OFF

void vControl() {

	initGlobalVariables();
	
	period = 500 ms;
		
	while(1) {

		ticks = xTaskGetTickCount();
		
		updateSensors();
		
		updateAlarms();
		
		controlActuators();
		
		vTaskDelayUntil(&ticks, period);
	}
}
\end{lstlisting}




%------------------------------------------------------------------------------
\subsection{EDU-CIAA-NXP}


%----------------------------------------------------------------------------------------
%	SECTION
%----------------------------------------------------------------------------------------
\section{Análisis del software}
 
La idea de esta sección es resaltar los problemas encontrados, los criterios utilizados y la justificación de las decisiones que se hayan tomado.

Se puede agregar código o pseudocódigo dentro de un entorno lstlisting con el siguiente código:

\begin{verbatim}
\begin{lstlisting}[caption= "un epígrafe descriptivo"]

	las líneas de código irían aquí...
	
\end{lstlisting}
\end{verbatim}

A modo de ejemplo:

\begin{lstlisting}[caption=Pseudocódigo del lazo principal de control.]  % Start your code-block

#define MAX_SENSOR_NUMBER 3
#define MAX_ALARM_NUMBER  6
#define MAX_ACTUATOR_NUMBER 6

uint32_t sensorValue[MAX_SENSOR_NUMBER];		
FunctionalState alarmControl[MAX_ALARM_NUMBER];	//ENABLE or DISABLE
state_t alarmState[MAX_ALARM_NUMBER];						//ON or OFF
state_t actuatorState[MAX_ACTUATOR_NUMBER];			//ON or OFF

void vControl() {

	initGlobalVariables();
	
	period = 500 ms;
		
	while(1) {

		ticks = xTaskGetTickCount();
		
		updateSensors();
		
		updateAlarms();
		
		controlActuators();
		
		vTaskDelayUntil(&ticks, period);
	}
}
\end{lstlisting}




%----------------------------------------------------------------------------------------
%	SECTION
%----------------------------------------------------------------------------------------
\section{Análisis del software}
 
La idea de esta sección es resaltar los problemas encontrados, los criterios utilizados y la justificación de las decisiones que se hayan tomado.

Se puede agregar código o pseudocódigo dentro de un entorno lstlisting con el siguiente código:

\begin{verbatim}
\begin{lstlisting}[caption= "un epígrafe descriptivo"]

	las líneas de código irían aquí...
	
\end{lstlisting}
\end{verbatim}

A modo de ejemplo:

\begin{lstlisting}[caption=Pseudocódigo del lazo principal de control.]  % Start your code-block

#define MAX_SENSOR_NUMBER 3
#define MAX_ALARM_NUMBER  6
#define MAX_ACTUATOR_NUMBER 6

uint32_t sensorValue[MAX_SENSOR_NUMBER];		
FunctionalState alarmControl[MAX_ALARM_NUMBER];	//ENABLE or DISABLE
state_t alarmState[MAX_ALARM_NUMBER];						//ON or OFF
state_t actuatorState[MAX_ACTUATOR_NUMBER];			//ON or OFF

void vControl() {

	initGlobalVariables();
	
	period = 500 ms;
		
	while(1) {

		ticks = xTaskGetTickCount();
		
		updateSensors();
		
		updateAlarms();
		
		controlActuators();
		
		vTaskDelayUntil(&ticks, period);
	}
}
\end{lstlisting}





%----------------------------------------------------------------------------------------
%	SECTION
%----------------------------------------------------------------------------------------
\section{Análisis del software}
 
La idea de esta sección es resaltar los problemas encontrados, los criterios utilizados y la justificación de las decisiones que se hayan tomado.

Se puede agregar código o pseudocódigo dentro de un entorno lstlisting con el siguiente código:

\begin{verbatim}
\begin{lstlisting}[caption= "un epígrafe descriptivo"]

	las líneas de código irían aquí...
	
\end{lstlisting}
\end{verbatim}

A modo de ejemplo:

\begin{lstlisting}[caption=Pseudocódigo del lazo principal de control.]  % Start your code-block

#define MAX_SENSOR_NUMBER 3
#define MAX_ALARM_NUMBER  6
#define MAX_ACTUATOR_NUMBER 6

uint32_t sensorValue[MAX_SENSOR_NUMBER];		
FunctionalState alarmControl[MAX_ALARM_NUMBER];	//ENABLE or DISABLE
state_t alarmState[MAX_ALARM_NUMBER];						//ON or OFF
state_t actuatorState[MAX_ACTUATOR_NUMBER];			//ON or OFF

void vControl() {

	initGlobalVariables();
	
	period = 500 ms;
		
	while(1) {

		ticks = xTaskGetTickCount();
		
		updateSensors();
		
		updateAlarms();
		
		controlActuators();
		
		vTaskDelayUntil(&ticks, period);
	}
}
\end{lstlisting}




%------------------------------------------------------------------------------
%	SECTION
%------------------------------------------------------------------------------
\section{Implementación del código C dependiende del hardware}
\label{sec:codeImplementation}

%------------------------------------------------------------------------------
%\subsection{}

En la implementación para las plataformas del proyecto CIAA, del código dependiente del hardware, que forma parte de la biblioteca sAPI, se utilizaron las bibliotecas de drivers provistas por los fabricantes. En particular, para los microcontroladores de la empresa NXP se utilizó LPCOpen, en la version 3.02 para el LPC4337 y versión 3.04 para el LPC54102. Para el microcontrolador EFM32HG322 de Siicon Labs se utilizó la biblioteca EMLIB versión 5.1.2 (todas en sus últimas versiones al momento de la realización de este trabajo).

Por otra parte, se reutilizó parte del código desarrollado en versiones anteriores de la biblioteca sAPI para las plataformas EDU-CIAA-NXP y CIAA-Z3R0 y parte del código de PicoFirmware \citep{PicoAPI} para la implementación en la plataforma PicoCIAA.

Una aplicación de embebidos típica que utiliza la biblioteca sAPI, se puede combinar con un sistema operativo de tiempo real, \emph{stracks} y \emph{middelware} resultando una arquitectura de capas de software típica en aplicaciones de sistemas embebidos como se ilustra en la figura \ref{fig:sapiCapas2}.

\begin{figure}[!htbp]
\begin{center}  % [width=14cm,height=8cm] [width=\textwidth]
\includegraphics*[width=10.4cm]{Figures/sapiCapas2.png}
\par\caption{Arquitectura de una aplicación que utiliza la biblioteca.}\label{fig:sapiCapas2}
\end{center}
\end{figure}


%----------------------------------------------------------------------------------------
%	SECTION
%----------------------------------------------------------------------------------------
\section{Análisis del software}
 
La idea de esta sección es resaltar los problemas encontrados, los criterios utilizados y la justificación de las decisiones que se hayan tomado.

Se puede agregar código o pseudocódigo dentro de un entorno lstlisting con el siguiente código:

\begin{verbatim}
\begin{lstlisting}[caption= "un epígrafe descriptivo"]

	las líneas de código irían aquí...
	
\end{lstlisting}
\end{verbatim}

A modo de ejemplo:

\begin{lstlisting}[caption=Pseudocódigo del lazo principal de control.]  % Start your code-block

#define MAX_SENSOR_NUMBER 3
#define MAX_ALARM_NUMBER  6
#define MAX_ACTUATOR_NUMBER 6

uint32_t sensorValue[MAX_SENSOR_NUMBER];		
FunctionalState alarmControl[MAX_ALARM_NUMBER];	//ENABLE or DISABLE
state_t alarmState[MAX_ALARM_NUMBER];						//ON or OFF
state_t actuatorState[MAX_ACTUATOR_NUMBER];			//ON or OFF

void vControl() {

	initGlobalVariables();
	
	period = 500 ms;
		
	while(1) {

		ticks = xTaskGetTickCount();
		
		updateSensors();
		
		updateAlarms();
		
		controlActuators();
		
		vTaskDelayUntil(&ticks, period);
	}
}
\end{lstlisting}





%------------------------------------------------------------------------------
%	SECTION
%------------------------------------------------------------------------------
\section{Documentación y difusión}
\label{sec:documentation}


%------------------------------------------------------------------------------
\subsection{Generación automática de manual de referencia de la API en base al modelo}



%------------------------------------------------------------------------------
\subsection{Tutoriales de instalación y uso}



%------------------------------------------------------------------------------
\subsection{Difusión a la comunidad del Proyecto CIAA y Embebidos32}


% Ensayos y Resultados
\chapter{Ensayos y Resultados}
\label{ChapterEnsayosYResultados}

\emph{Pruebas funcionales del hardware: La idea de esta sección es explicar cómo se hicieron los ensayos, qué resultados se obtuvieron y analizarlos.}

%------------------------------------------------------------------------------
%	SECTION
%------------------------------------------------------------------------------
\section{Testeo Unitario}
\label{sec:unitTest}


%------------------------------------------------------------------------------
%	SECTION
%------------------------------------------------------------------------------
\section{Banco de pruebas de hardware}
\label{sec:testBench}


%------------------------------------------------------------------------------
%	SECTION
%------------------------------------------------------------------------------
\section{Integración contínua}
\label{sec:ci}


%------------------------------------------------------------------------------
%	SECTION
%------------------------------------------------------------------------------
\section{Utilización de la biblioteca para la enseñanza de programación de Sistemas Embebidos}
\label{sec:teach}


----------------------------------------------

CASOS DE USO (Paper sAPI)

Desde la primer versión de la biblioteca sAPI realizada en el marco del proyecto de Java[] en 2015 se ha utilizado como ejemplo de capa de abstracción de hardware en las asignaturas dictadas por el autor. Estas asignaturas son: el  curso de posgrado “Programación de microproce- sadores”[] de la Carrera Especialización en Sistemas Embebidos (CESE[]) de la FI-UBA. Donde el autor se ha desempeñado como docente a cargo del curso durante tres ediciones, y la asignatura “Sistemas Digitales” de la carrera Ingeniería en Automatización y Control Industrial (IACI[]) de la Universidad Nacional de Quilmes (UNQ) donde el autor se desempeña en la actualidad como instructor.

A partir mazo de 2016 se decide utilizar la biblioteca como en el marco de los Cursos Abiertos de Programación de Sistemas Embebidos (CAPSE[]) organizados por la ACSE[]. De esta forma se ha extendido la biblioteca de forma considerable para explicar la utilización de todos los periféricos típicos de microcontroladores con excelentes resultados en cuanto a aprendizaje por parte de los alumnos tanto de niveles avanzados como quienes dan sus primeros pasos en el aprendizaje de programación de microcontroladores.
Además, la biblioteca sAPI se puso a disposición de cualquier persona ya que se encuentra publicada de forma libre y gratuita por internet bajo una licencia BSD modificada[] en el sitio de github del autor[] y ha sido adoptada por una gran cantidad de usuarios. 

Finalmente, en diciembre de 2016 se decide utilizar la biblioteca como biblioteca estándar para las plataformas del Proyecto CIAA. Esto llevó a una profunda revisión y mejora de la misma y todavía se está trabajando en la actualidad.


----------------------------------------------


CASOS DE USO (paper tools)

Estas herramientas se han utilizado en múltiples, entre
ellos, el curso "Sistemas digitales" de la carrera Ingeniería en
Automatización y Control Industrial en la UNQ, donde el
autor actualmente se desempeña como Profesor Instructor
desde 2014. El plan de estudio de este curso incluye
programación avanzada en lenguaje C para
microcontroladores, con temas tales como modularización de
código, máquinas de estado finito, sistemas operativos en
tiempo real que usan programación cooperativa y apropiativa
en microcontroladores; el curso de posgrado "Programación
de microprocesadores" dentro de la CESE de FI-UBA donde
el autor ha sido profesor a cargo de tres ediciones. Este curso
incluye: programación básica de microcontroladores en
lenguaje C, modularización, máquinas de estados finitos.
También en cursos organizados por ACSE []. En estos
cursos las herramientas junto con la secuencia didáctica
propuesta fueron intensamente utilizadas, mejorando esta
secuencia entre el autor y Pablo Gómez. En los cursos de
ACSE se distinguen dos grupos, "Cursos Abiertos de
Programación de Sistemas Embebidos" (CAPSE), orientados a
cualquier persona interesada en aprender la programación de
Sistemas Embebidos; y cursos impartidos al "Instituto
Nacional de Escuelas Técnicas" (INET) [], enfocados en la
capacitación de docentes de Escuelas Técnicas Secundarias,
para posteriormente poder retransmitir a sus alumnos en todo
el país. En ambos casos, se han observado muy buenos
resultados de aprendizaje. En el caso de los cursos CAPSE, se
impartieron cuatro cohortes entre 2016 y 2017; en la primera
cohorte, el 46\% de los estudiantes completaron todos los
niveles con un promedio de 67\% de índice de aprobación por
nivel; mientras que en la cuarta cohorte, el 73\% completó
todos los niveles, con una tasa de aprobación promedio de
81\% por nivel. En el caso de los cursos INET, con dos
cohortes en 2017, el 68\% de los estudiantes completaron todos
los niveles, con una tasa promedio de aprobación del 74\% por
nivel. Además, después de completar sus estudios, algunos de
estos estudiantes se inscribieron posteriormente en la CESE
FI-UBA para continuar profundizando su aprendizaje.
Además, se han realizado varios talleres en escuelas
secundarias y universidades de todo el país. Vale la pena
mencionar los múltiples tutoriales y workshop en el "Simposio
Argentino de Sistemas Embebidos" (SASE) [].
Todas estas herramientas son de código abierto, publicadas de
forma gratuita en la cuenta web de github del Proyecto CIAA
y han sido adoptadas por un número importante de usuarios y
docentes.

% Conclusiones
\chapter{Conclusiones} % Main chapter title
\label{ChapterConclusiones} % Change X to a consecutive number; for referencing
                            % this chapter elsewhere, use \ref{ChapterX}

%----------------------------------------------------------------------------------------
%	SECTION
%----------------------------------------------------------------------------------------

\section{Trabajo realizado}

\emph{La idea de esta sección es resaltar cuáles son los principales aportes del trabajo realizado y cómo se podría continuar. Debe ser especialmente breve y concisa. Es buena idea usar un listado para enumerar los logros obtenidos.}

%----------------------------------------------------------------------------------------
%	SECTION
%----------------------------------------------------------------------------------------
\section{Próximos pasos}

\emph{Acá se indica cómo se podría continuar el trabajo más adelante.}
 

%------------------------------------------------------------------------------
%	CONTENIDO DE LA MEMORIA  - APÉNDICES
%------------------------------------------------------------------------------

% indicativo para indicarle a LaTeX los siguientes "capítulos" son apéndices
\appendix 

% Incluir los apéndices de la memoria como archivos separadas desde la carpeta
% Appendices
% Descomentar las líneas a medida que se escriben los apéndices

%% Appendix A

\chapter{Appendix Title Here} % Main appendix title

\label{AppendixA} % For referencing this appendix elsewhere, use \ref{AppendixA}

Write your Appendix content here.
%\include{Appendices/AppendixB}
%\include{Appendices/AppendixC}

%------------------------------------------------------------------------------
%	BIBLIOGRAPHY
%------------------------------------------------------------------------------

\Urlmuskip=0mu plus 1mu\relax
\raggedright
\printbibliography[heading=bibintoc]

%------------------------------------------------------------------------------

\end{document}  
